% !TeX program = xelatex
% !TeX encoding = UTF-8
\documentclass{article}
\usepackage{ctex}
\usepackage{geometry}
\geometry{a4paper, left=2.5cm, right=2.5cm, top=2.5cm, bottom=2.5cm}

\usepackage{amsmath, amssymb}
\usepackage{array}
\usepackage{booktabs}
\usepackage{listings}
\usepackage{algorithm, algorithmic}
\floatname{algorithm}{算法}
\renewcommand{\algorithmicrequire}{\textbf{输入:}}
\renewcommand{\algorithmicensure}{\textbf{输出:}}
\usepackage{xcolor}
\usepackage{hyperref}
\hypersetup{colorlinks, citecolor=black, filecolor=black, linkcolor=black, urlcolor=black}

% 代码样式
\lstset{
    language=Python,
    keywordstyle=\color{blue},
    commentstyle=\color{green!50!black},
    stringstyle=\color{red},
    numbers=left,
    numberstyle=\tiny,
    basicstyle=\ttfamily\small,
    frame=single,
    breaklines=true,
    extendedchars=false
}

\title{第九次上机作业：数值积分方法比较}
\author{学号：241830137\quad 姓名：朱子健}
\date{\today}

\begin{document}

\maketitle

\section{问题描述}

数学上可以证明
\begin{equation}
    \int_{0}^{1} \frac{4}{1 + x^{2}} \, \mathrm{d}x = \pi .
\end{equation}
本题分别使用复合梯形公式和复合 Simpson 公式计算 $\pi$ 的近似值。 选择不同的 $h$（或等分数 $n$），对每种求积公式，试将误差刻画成 $h$ 的函数，并比较两种方法的精度。 特别关注是否存在某个 $h$，当低于这个值之后再继续减小 $h$，计算不再有所改进，并解释其原因。

\section{算法原理}

\subsection{复合梯形公式}
将积分区间 $[a,b]$ 作 $n$ 等分，步长 $h = (b-a)/n$，节点 $x_i = a + i h$，$i=0,1,\dots,n$。 复合梯形求积公式为
\begin{equation}
    T_n(f) = \frac{h}{2}\left[ f(a) + f(b) + 2\sum_{i=1}^{n-1} f(a + i h) \right].
\end{equation}
其离散误差为 $O(h^2)$。

\begin{algorithm}[H]
\caption{复合梯形公式}
\begin{algorithmic}[1]
\REQUIRE 区间端点 $a, b$，等分数 $n$，被积函数 $f(x)$
\ENSURE 积分近似值 $T$
\STATE $h \gets (b - a) / n$
\STATE $s \gets 0.5 \cdot f(a) + 0.5 \cdot f(b)$
\FOR {$i \gets 1$ \TO $n-1$}
    \STATE $x \gets a + i \cdot h$
    \STATE $s \gets s + f(x)$
\ENDFOR
\STATE $T \gets h \cdot s$
\RETURN $T$
\end{algorithmic}
\end{algorithm}

\subsection{复合 Simpson 公式}
将区间 $[a,b]$ 作 $n$ 等分，其中 $n=2m$ 为偶数，步长 $h = (b-a)/n$。 在每两个子区间上应用 Simpson 公式，得到复合 Simpson 求积公式
\begin{equation}
    S_m(f) = \frac{h}{3}\left[ f(a) + f(b) + 4\sum_{i=1}^{m} f(a + (2i-1)h) + 2\sum_{i=1}^{m-1} f(a + 2i h) \right].
\end{equation}
其离散误差为 $O(h^4)$。

\begin{algorithm}[H]
\caption{复合 Simpson 公式}
\begin{algorithmic}[1]
\REQUIRE 区间端点 $a, b$，等分数 $n$（$n$ 为偶数），被积函数 $f(x)$
\ENSURE 积分近似值 $S$
\IF {$n$ 为奇数}
    \STATE $n \gets n + 1$
\ENDIF
\STATE $h \gets (b - a) / n$
\STATE $s \gets f(a) + f(b)$
\FOR {$i \gets 1$ \TO $n-1$}
    \STATE $x \gets a + i \cdot h$
    \IF {$i$ 为奇数}
        \STATE $s \gets s + 4 \cdot f(x)$
    \ELSE
        \STATE $s \gets s + 2 \cdot f(x)$
    \ENDIF
\ENDFOR
\STATE $S \gets h \cdot s / 3$
\RETURN $S$
\end{algorithmic}
\end{algorithm}

\subsection{误差停止改进的原因}
随着 $h \to 0$，离散误差按 $O(h^2)$ 或 $O(h^4)$ 减小，但浮点舍入误差随节点数增加而累积。 当 $h$ 极小（$n$ 极大）时，舍入误差主导总误差，导致误差不再减小甚至回升。 在单精度（float32）下这一现象比双精度出现得更早且更明显。

\section{结果分析}

为了清晰地显示舍入误差的影响，程序采用单精度浮点数（\texttt{np.float32}）进行计算，并将 $n$ 取至非常大。 精确值取为 $\pi$ 的双精度值。

\begin{table}[htbp]
\centering
\caption{复合梯形和复合 Simpson 公式在 float32 下的误差变化}
\label{tab:errors}
\begin{tabular}{c c c c c c}
\toprule
$n$ & $h$ & 梯形误差 & 改进？& Simpson 误差 & 改进？ \\
\midrule
2       & 5.00e-01 & 2.14e-02  & --- & 1.14e-03  & --- \\
4       & 2.50e-01 & 5.35e-03  & 是 & 7.41e-05  & 是 \\
8       & 1.25e-01 & 1.33e-03  & 是 & 4.64e-06  & 是 \\
16      & 6.25e-02 & 3.31e-04  & 是 & 1.91e-06  & 是 \\
32      & 3.13e-02 & 8.47e-05  & 是 & 5.96e-07  & 是 \\
64      & 1.56e-02 & 2.87e-05  & 是 & 1.79e-06  & 否 \\
128     & 7.81e-03 & 1.70e-05  & 是 & 9.54e-07  & 是 \\
256     & 3.91e-03 & 2.54e-05  & 否 & 4.17e-06  & 否 \\
512     & 1.95e-03 & 1.14e-04  & 否 & 4.17e-06  & 否 \\
1024    & 9.77e-04 & 9.01e-05  & 是 & 4.17e-06  & 否 \\
2048    & 4.88e-04 & 6.13e-04  & 否 & 4.77e-06  & 否 \\
4096    & 2.44e-04 & 4.96e-04  & 是 & 2.41e-05  & 否 \\
8192    & 1.22e-04 & 1.83e-03  & 否 & 5.96e-05  & 否 \\
16384   & 6.10e-05 & 4.97e-04  & 是 & 5.72e-05  & 是 \\
32768   & 3.05e-05 & 6.65e-03  & 否 & 3.25e-04  & 否 \\
65536   & 1.53e-05 & 3.33e-03  & 是 & 1.16e-03  & 否 \\
\bottomrule
\end{tabular}
\end{table}

从表 \ref{tab:errors} 可以看出：
\begin{itemize}
    \item 在 $n$ 较小（$h$ 较大）时，梯形公式误差约为 $O(h^2)$，Simpson 公式误差约为 $O(h^4)$，后者精度明显更高。
    \item 当 $n$ 增大到一定值后，两种公式的误差均停止稳步下降，出现波动甚至回升。 梯形公式在 $n=32$ 左右误差最小，Simpson 公式在 $n=32\sim 128$ 左右误差最小。
    \item 此后继续增大 $n$（减小 $h$），舍入误差逐渐主导，数值结果不再改进。 由于使用 float32 单精度，表示精度有限（机器精度约 $1.2\times10^{-7}$），大量函数值相加引入的舍入噪声完全淹没了离散误差的减小。
\end{itemize}

\section{结论}

\begin{itemize}
    \item 复合 Simpson 公式在同等步长下具有远优于复合梯形公式的精度，其误差以 $O(h^4)$ 收敛，而梯形公式仅为 $O(h^2)$。
    \item 无论是复合梯形还是复合 Simpson 公式，都存在一个最优的 $h$（或 $n$）。 当步长小于该值时，由于计算机浮点运算的舍入误差累积，近似值的总误差不再下降，甚至增大。
    \item 在 float32 单精度下，该现象非常显著；若改用 double 双精度，最优 $h$ 可以小得多，但最终也会受到同样的限制。
\end{itemize}

\section{附录：程序代码}

\subsection{复合梯形与复合 Simpson 公式（float32）}
\begin{lstlisting}[language=Python]
import numpy as np
def f32(x):
    """单精度被积函数"""
    return np.float32(4.0) / (np.float32(1.0) + np.float32(x)**2)

def composite_trapezoidal_float32(a, b, n):
    """复合梯形公式，float32单精度"""
    h = np.float32((b - a) / n)
    s = np.float32(0.5) * f32(a) + np.float32(0.5) * f32(b)
    for i in range(1, n):
        s += f32(a + np.float32(i) * h)
    return float(h * s)

def composite_simpson_float32(a, b, n):
    """复合Simpson公式，float32单精度"""
    if n % 2 != 0:
        n += 1
    h = np.float32((b - a) / n)
    s = f32(a) + f32(b)
    for i in range(1, n, 2):
        s += np.float32(4.0) * f32(a + np.float32(i) * h)
    for i in range(2, n, 2):
        s += np.float32(2.0) * f32(a + np.float32(i) * h)
    return float(h / np.float32(3.0) * s)


a, b = 0, 1
exact = np.pi

print("使用float32单精度，误差不再减小的现象更显著：\n")
print(f"{'n':<10} {'h':<14} {'梯形误差':<16} {'有改进?':<10} {'Simpson误差':<16} {'有改进?':<10}")
print("-" * 76)

prev_T = None
prev_S = None

for n in [2, 4, 8, 16, 32, 64, 128, 256, 512, 1024, 2048, 4096, 8192, 16384, 32768, 65536]:
    h = (b - a) / n
    err_T = abs(composite_trapezoidal_float32(a, b, n) - exact)
    err_S = abs(composite_simpson_float32(a, b, n) - exact)

    if prev_T is None:
        T_better, S_better = "—", "—"
    else:
        T_better = "是" if err_T < prev_T else "否"
        S_better = "是" if err_S < prev_S else "否"

    print(f"{n:<10} {h:<14.6e} {err_T:<16.6e} {T_better:<10} {err_S:<16.6e} {S_better:<10}")

    prev_T = err_T
    prev_S = err_S
\end{lstlisting}
\end{document}